\section{Implementacja systemu}

W poprzednim rozdziale skupiono się na projekcie systemu, przedstawiono jego główne elementy oraz dane, które będą pomiędzy nimi przepływać. Niniejszy rozdział opisuje implementację zaprezentowanego planu. Przedstawia, w jaki sposób w kodzie źródłowym zostały zrealizowane poszczególne elementy niezbędne do prawidłowego działania aplikacji.

Zgodnie z tym, co przedstawiono w projekcie, system zaimplementowano jako aplikację sieciową składającą się z trzech głównych części: aplikacji klienckiej, serwerowej oraz modułu analizy wideo. Każda z tych części ma na celu wykonywanie szczegółowych zadań opisanych w części projektowej.

Opis implementacji uporządkowano zgodnie z przebiegiem działania systemu. W pierwszej kolejności przedstawiono środowisko implementacyjne i organizację kodu źródłowego. Następnie opisano moduł odpowiedzialny za analizę nagrania wideo, ponieważ stanowi on najważniejszą część systemu z punktu widzenia celu pracy. W dalszej części rozdziału omówiono implementację aplikacji serwerowej, interfejsu użytkownika aplikacji klienckiej oraz sposób integracji wszystkich komponentów.

\subsection{Środowisko implementacyjne i organizacja kodu źródłowego}

Kod źródłowy systemu został umieszczony w repozytorium prowadzonym z wykorzystaniem systemu kontroli wersji Git. Repozytorium zawiera kompletną strukturę projektu, pliki konfiguracyjne oraz historię zmian wprowadzanych podczas implementacji. W katalogu głównym projektu znajdują się trzy główne katalogi: \texttt{backend}, \texttt{frontend} oraz \texttt{ml\_model}. Oprócz nich umieszczono tam również między innymi plik \texttt{requirements.txt}, zawierający zależności części serwerowej i modułu analizy, plik \texttt{README.md} oraz plik \texttt{.gitignore}, określający elementy pomijane przez system kontroli wersji.

Taki układ repozytorium porządkuje kod według odpowiedzialności poszczególnych części systemu. Katalog \texttt{backend} zawiera aplikację serwerową, katalog \texttt{frontend} zawiera aplikację kliencką, natomiast katalog \texttt{ml\_model} zawiera kod odpowiedzialny za analizę nagrania wideo. Dokonany podział pozwala zaakcentować separację pomiędzy poszczególnymi częściami aplikacji, zwiększając czytelność rozwiązania znajdującego się w repozytorium~\cite{microsoft_arch_principles}.
\newpage
Uproszczoną strukturę katalogów projektu przedstawiono na rysunku~\ref{fig:struktura-repozytorium}.

\begin{figure}[H]
\centering
\begin{minipage}{0.85\textwidth}
\dirtree{%
.1 Dead-Lift-Evaluation/.
.2 backend/.
.3 api/.
.3 config/.
.3 manage.py.
.2 frontend/.
.3 public/.
.3 src/.
.4 components/.
.4 config/.
.4 services/.
.4 utils/.
.4 App.jsx.
.4 main.jsx.
.3 package.json.
.2 ml\_model/.
.3 analysis/.
.3 models/.
.3 tests/.
.3 utils/.
.3 constants.py.
.3 main.py.
.2 .gitignore.
.2 README.md.
.2 requirements.txt.
}
\end{minipage}
\caption{Uproszczona struktura katalogów repozytorium projektu. Źródło: opracowanie własne.}
\label{fig:struktura-repozytorium}
\end{figure}


Część serwerowa została zaimplementowana w języku Python z użyciem szkieletu Django i udostępnia interfejs API w architekturze REST. Jej kod znajduje się w katalogu \texttt{backend}, zawierającym między innymi plik \texttt{manage.py}, katalog konfiguracyjny \texttt{config} oraz aplikację \texttt{api}. Część ta odpowiada za obsługę żądań, zapis danych związanych z analizą oraz komunikację z modułem analizy nagrania.

Część kliencka została zaimplementowana w języku JavaScript z użyciem biblioteki React oraz narzędzia Vite. Jej kod znajduje się w katalogu \texttt{frontend}. W katalogu \texttt{src} wydzielono między innymi katalogi \texttt{components}, \texttt{services}, \texttt{config} oraz \texttt{utils}, odpowiadające za poszczególne elementy interfejsu, komunikację z backendem, konfigurację oraz funkcje pomocnicze.

Moduł analizy nagrania wideo został zaimplementowany w języku Python i znajduje się w katalogu \texttt{ml\_model}. W jego strukturze wydzielono między innymi katalog \texttt{analysis} zawierający kod właściwej analizy, katalog \texttt{models} przechowujący pliki wykorzystywanych modeli, katalog \texttt{tests} oraz katalog \texttt{utils} z funkcjami pomocniczymi. W katalogu głównym modułu znajduje się również plik \texttt{constants.py}, zawierający stałe wykorzystywane podczas analizy.

Zależności części napisanej w języku programowania Python zostały określone w pliku \texttt{requirements.txt}. W projekcie wykorzystano między innymi Django, Django REST Framework, MediaPipe, OpenCV, NumPy oraz Matplotlib. Biblioteki te odpowiadają kolejno za implementację aplikacji serwerowej, udostępnienie interfejsu REST API, detekcję punktów kluczowych sylwetki, przetwarzanie obrazu, wykonywanie obliczeń numerycznych oraz pomocniczą obsługę danych i wizualizacji. Zależności części klienckiej oraz skrypty uruchomieniowe zostały określone w pliku \texttt{package.json}. Plik ten wskazuje użycie biblioteki React, React DOM, narzędzia Vite oraz narzędzi pomocniczych związanych z kontrolą jakości kodu JavaScript.

W repozytorium zastosowano również plik \texttt{.gitignore}, którego zadaniem jest określenie plików i katalogów pomijanych przez system kontroli wersji~\cite{gitignore_docs}. W projekcie wykluczono między innymi pliki generowane automatycznie podczas działania Pythona, lokalne środowiska wirtualne, pliki tymczasowe, lokalną bazę danych oraz katalogi i pliki związane z przesyłanymi lub testowymi materiałami wideo. Dzięki temu repozytorium jest utrzymywane w czystej postaci, zawiera kod źródłowy i pliki konfiguracyjne, ale pomija pliki wynikowe oraz należące do lokalnego środowiska programistycznego.

\subsection{Implementacja modułu analizy nagrania wideo}

Najważniejszą częścią systemu z punktu widzenia celu pracy jest moduł odpowiedzialny za analizę nagrania wideo. Jego zadaniem jest przetworzenie pliku przesłanego przez użytkownika, wykrycie sylwetki ćwiczącego, wyznaczenie podstawowych parametrów ruchu oraz przygotowanie danych potrzebnych do oceny techniki wykonania martwego ciągu. Kod tego modułu został umieszczony w katalogu \texttt{ml\_model}, natomiast główne funkcje związane z analizą znajdują się w katalogu \texttt{analysis}.

Centralnym elementem modułu jest klasa \texttt{DeadliftAnalyzer}. Nie odpowiada ona bezpośrednio za wszystkie obliczenia, lecz koordynuje kolejne etapy przetwarzania nagrania. Takie rozwiązanie pozwala zachować czytelny przepływ analizy i oddzielić od siebie fragmenty kodu odpowiedzialne za inne zadania. W konstruktorze klasy tworzony jest obiekt estymatora pozy, jeżeli nie został on przekazany z zewnątrz. Dzięki temu klasa może zostać uruchomiona z domyślnym sposobem detekcji punktów kluczowych, ale jednocześnie pozostaje możliwe przekazanie innego obiektu spełniającego tę samą rolę.

Funkcjonalność klasy \texttt{DeadliftAnalyzer} przedstawiono we fragmencie kodu~\ref{lst:deadlift-analyzer}. Na wycinku kodu widać, że nie implementuje ona bezpośrednio wszystkich szczegółowych obliczeń, lecz wywołuje kolejne funkcje odpowiedzialne za poszczególne etapy analizy, zachowując abstrakcję zgodną ze standardami paradygmatu programowania obiektowego.

\begin{lstlisting}[
language=Python,
caption={Klasa koordynująca proces analizy martwego ciągu. Źródło: opracowanie własne.},
label={lst:deadlift-analyzer},
basicstyle=\small\ttfamily,
breaklines=true,
breakatwhitespace=true,
columns=fullflexible
]
class DeadliftAnalyzer:
    def __init__(self, estimator=None):
        self.estimator = estimator or PoseEstimator()

    def analyze(self, video_path: str) -> dict:
        timeline = analyze_video(video_path, self.estimator)

        repetitions = detect_repetitions(timeline)
        timeline = assign_repetitions_to_timeline(timeline, repetitions)

        summaries = summarize_repetitions(timeline, repetitions)

        evaluations = evaluate_repetitions(summaries)
        consistency_evaluation = evaluate_repetition_consistency(summaries)

        return {
            "repetitions_count": len(repetitions),
            "timeline": timeline,
            "repetitions": repetitions,
            "summaries": summaries,
            "evaluations": evaluations,
            "consistency_evaluation": consistency_evaluation,
        }
\end{lstlisting}

Proces analizy rozpoczyna się od wywołania metody \texttt{analyze}, która przyjmuje ścieżkę do pliku wideo. Ścieżka ta nie jest ustawiona na stałe w klasie analizującej, lecz przekazywana z zewnątrz. Ma to znaczenie przy integracji z aplikacją serwerową, ponieważ nagranie pochodzi wtedy z pliku przesłanego przez użytkownika. Dzięki takiemu podejściu ten sam moduł może zostać wykorzystany zarówno z poziomu aplikacji serwerowej, jak i podczas lokalnych testów prowadzonych bez udziału interfejsu użytkownika.

Metoda \texttt{analyze} koordynuje kolejne etapy przetwarzania nagrania. W pierwszej kolejności tworzona jest oś czasu zawierająca dane uzyskane z analizowanych klatek. Następnie wykrywane są powtórzenia i przypisywane są im fazy ruchu. Na ich podstawie tworzone są podsumowania pojedynczych powtórzeń, które następnie wykorzystywane są do oceny techniki oraz spójności całej serii.

Wynikiem działania metody jest słownik zawierający dane powstałe na kolejnych etapach analizy, między innymi oś czasu, listę wykrytych powtórzeń, ich podsumowania i oceny. Szczegółowy sposób realizacji poszczególnych etapów przedstawiono w kolejnych podrozdziałach.

\subsection{Detekcja punktów kluczowych i wyznaczanie metryk biomechanicznych}

Pierwszym etapem właściwej analizy nagrania jest przetworzenie filmu na pojedyncze klatki oraz wykrycie na nich punktów kluczowych sylwetki. W projekcie za iterowanie po klatkach odpowiada funkcja \texttt{iter\_video\_frames}, znajdująca się w katalogu \texttt{utils}. Funkcja ta otwiera plik wideo za pomocą biblioteki OpenCV, a następnie zwraca kolejne klatki wraz z numerem klatki oraz czasem nagrania wyrażonym w sekundach. Analiza nie jest wykonywana dla każdej klatki filmu, lecz z uwzględnieniem parametru \texttt{frame\_stride}. Wartość domyślna tego parametru została określona w pliku \texttt{constants.py} jako stała \texttt{VIDEO\_FRAME\_STRIDE}. Dzięki temu możliwe jest ograniczenie liczby analizowanych klatek przy zachowaniu informacji o przebiegu ruchu w czasie.

Do wykrywania punktów kluczowych sylwetki wykorzystano klasę \texttt{PoseEstimator}. Klasa ta korzysta z gotowego modelu MediaPipe Pose Landmarker, którego ścieżka została określona w stałej \texttt{MODEL\_PATH}. W konstruktorze klasy tworzony jest obiekt detektora pozy, skonfigurowany do pracy w trybie analizy pojedynczego obrazu. W przypadku analizy filmu oznacza to, że każda wybrana klatka nagrania jest traktowana jako osobny obraz przekazywany do modelu estymacji pozy.

Metoda \texttt{detect\_pose\_from\_frame} przyjmuje klatkę filmu odczytaną przez OpenCV. Ponieważ OpenCV przechowuje obraz w formacie BGR, klatka jest najpierw konwertowana do formatu RGB. Następnie tworzony jest obiekt obrazu zgodny z wymaganiami biblioteki MediaPipe i wykonywana jest detekcja punktów kluczowych. Jeżeli model nie wykryje sylwetki, metoda zwraca wartość \texttt{None}. Taka sytuacja jest później zapisywana w osi czasu analizy jako klatka, dla której nie wykryto pozy ćwiczącego.

W przypadku pomyślnie zakończonej detekcji wynik zwrócony zostaje w postaci słownika. Do każdego punktu kluczowego przyporządkowany jest stały indeks, a w wartości do niego przypisanej znajdują się współrzędne \texttt{x}, \texttt{y}, \texttt{z} oraz wartość \texttt{visibility}. Wartość \texttt{visibility} określa przewidywaną widoczność danego punktu na obrazie. Współrzędna \texttt{z} nie jest wykorzystywana w dalszych obliczeniach kątów, ponieważ analiza opiera się na współrzędnych dwuwymiarowych.

Kolejna część analizy wymaga odwoływania się do konkretnych stawów, takich jak bark, biodro, kolano, kostka czy nadgarstek. Z tego powodu w pliku \texttt{constants.py} zdefiniowano mapowanie nazw punktów na identyfikatory wykorzystywane przez model MediaPipe. Dzięki temu w pozostałych częściach programu można posługiwać się nazwami takimi jak \texttt{left\_hip}, \texttt{right\_knee} czy \texttt{left\_shoulder}, zamiast bezpośrednio odwoływać się do numerów punktów.

Po wykryciu punktów kluczowych wybierana jest strona ciała wykorzystywana do dalszej analizy. Dla punktów objętych analizą, znajdujących się po obu stronach sylwetki, obliczany jest średni wynik widoczności. Następnie wybierana jest strona o wyższej średniej wartości \texttt{visibility}. Przy bocznym i statycznym ustawieniu kamery jedna strona sylwetki jest zwykle lepiej widoczna. Wybór tej strony ogranicza wpływ punktów częściowo zasłoniętych na dalsze obliczenia. W przypadku remisu domyślnie wybierana jest strona lewa.

Następnym krokiem jest obliczenie kątów pomiędzy wybranymi punktami sylwetki. Za tę część obliczeń odpowiada plik \texttt{angle\_calculation.py}. Funkcja zawarta w pliku o nazwie \texttt{calculate\_angle\_from\_points} oblicza kąt utworzony przez trzy punkty na płaszczyźnie dwuwymiarowej. Punkt środkowy traktowany jest jako wierzchołek kąta. Obliczenie wykonywane jest na podstawie współrzędnych \texttt{x} oraz \texttt{y}. Wyznaczane są dwa wektory, ich długości, iloczyn skalarny oraz wartość cosinusa kąta. Następnie wynik jest przeliczany z radianów na stopnie. W kodzie uwzględniono również zabezpieczenie przed dzieleniem przez zero oraz przed przekazaniem do funkcji odwrotnego cosinusa wartości spoza zakresu od \texttt{-1} do \texttt{1}.

Na podstawie funkcji obliczającej kąt dla trzech punktów zbudowano funkcję o nazwie \texttt{calculate\_angle\_from\_joints}. Przyjmuje ona nazwy trzech stawów, pobiera odpowiadające im współrzędne i przekazuje je do funkcji liczącej kąt. Takie rozwiązanie pozwala definiować kąty w sposób bardziej czytelny, ponieważ w kodzie można wskazać nazwy stawów zamiast ręcznie pobierać współrzędne poszczególnych punktów.

Główne metryki biomechaniczne dla wybranej strony ciała są obliczane przez funkcję \texttt{calculate\_single\_side\_biomechanical\_metrics}. W pierwszej kolejności wyznaczane są kąty opisane w stałej \texttt{SINGLE\_SIDE\_ANGLE\_DEFINITIONS}. W obecnej wersji systemu są to kąty: \texttt{knee\_angle}, \texttt{hip\_angle} oraz \texttt{shoulder\_angle}. Definicje te określają, które trzy punkty sylwetki mają zostać użyte do obliczenia danego kąta. Przykładowo kąt kolana wyznaczany jest na podstawie biodra, kolana i kostki po wybranej stronie ciała.

Oprócz kątów stawowych obliczany jest również kąt pochylenia tułowia, zapisany jako \texttt{back\_inclination\_angle}. Metryka ta jest dodawana do słownika wynikowego razem z pozostałymi kątami. W przypadku gdy dla danego kąta lub metryki nie uda się wykonać obliczenia, do wyniku przypisywana jest wartość \texttt{None}. Dzięki temu system może kontynuować analizę kolejnych klatek, nawet jeżeli dla jednej z nich część danych nie została poprawnie wyznaczona.

W analizie uwzględniono również pozycję środka gryfu sztangi. Nie jest ona wykrywana jako osobny obiekt na obrazie, lecz wyznaczana na podstawie przybliżonego środka pomiędzy nadgarstkami. Jest to jedyna metryka, do której wyznaczenia wykorzystywane są punkty z obu stron ciała. Za wyznaczenie tej pozycji odpowiada funkcja \texttt{calculate\_bar\_position}, która zwraca środek pomiędzy lewym i prawym nadgarstkiem. Przybliżenie to wynika z założenia, że podczas martwego ciągu środek pomiędzy nadgarstkami znajduje się w pobliżu środka gryfu.

Wszystkie opisane elementy są łączone w funkcji \texttt{analyze\_video}. Dla każdej analizowanej klatki funkcja ta próbuje wykryć punkty kluczowe sylwetki, wybiera lepiej widoczną stronę ciała, oblicza metryki biomechaniczne oraz wyznacza przybliżoną pozycję sztangi. Wynikiem działania funkcji jest oś czasu analizy, w której każda przetworzona klatka posiada przypisane informacje o detekcji pozy, wybranej stronie ciała, obliczonych metrykach, położeniu sztangi oraz punktach kluczowych sylwetki. Tak przygotowane dane stanowią podstawę dla dalszych etapów systemu.

W celu ułatwienia interpretacji danych zapisywanych w osi czasu analizy, w tabeli~\ref{tab:znaczenie-parametrow-osi-czasu} przedstawiono znaczenie najważniejszych parametrów wykorzystywanych w dalszej części programu.

\begin{table}[H]
\centering
\caption{Znaczenie parametrów zapisywanych w osi czasu analizy. Źródło: opracowanie własne.}
\label{tab:znaczenie-parametrow-osi-czasu}
\vspace{0.3cm}
\small
\begin{tabular}{p{4cm}p{10cm}}
\toprule
\textbf{Parametr} & \textbf{Znaczenie} \\
\midrule
\texttt{Klatka} & Numer klatki nagrania analizowanej przez system. \\

\texttt{Czas [s]} & Czas odpowiadający danej klatce, wyrażony w sekundach od początku nagrania. \\

\texttt{Poza} & Informacja określająca, czy dla danej klatki wykryto sylwetkę ćwiczącego. \\

\texttt{Strona} & Strona ciała wybrana do dalszej analizy na podstawie widoczności punktów kluczowych. \\

\texttt{knee\_angle} & Kąt w stawie kolanowym, liczony na podstawie położenia biodra, kolana i kostki po wybranej stronie ciała. \\

\texttt{hip\_angle} & Kąt w stawie biodrowym, liczony na podstawie położenia barku, biodra i kolana po wybranej stronie ciała. \\

\texttt{shoulder\_angle} & Kąt w stawie barkowym, liczony na podstawie położenia biodra, barku i nadgarstka po wybranej stronie ciała. \\

\texttt{back\_inclination} & Kąt pochylenia tułowia względem pionu, liczony na podstawie położenia barku i biodra po wybranej stronie ciała. \\

\texttt{bar\_x} & Przybliżona pozioma współrzędna środka gryfu sztangi, wyznaczana na podstawie położenia nadgarstków. \\

\texttt{bar\_y} & Przybliżona pionowa współrzędna środka gryfu sztangi, wyznaczana na podstawie położenia nadgarstków. \\

\texttt{bar\_visibility} & Średnia wartość widoczności lewego i prawego nadgarstka, używana pomocniczo przy przybliżaniu pozycji środka gryfu sztangi. \\
\bottomrule
\end{tabular}
\end{table}


Przykładowe wpisy osi czasu analizy, utworzone po przetworzeniu kolejnych klatek nagrania, przedstawiono w tabeli~\ref{tab:przykladowe-wpisy-osi-czasu}.

\begin{table}[H]
\centering
\caption{Przykładowe wpisy osi czasu analizy dla kolejnych klatek nagrania, pochodzące ze środkowej fazy ruchu. Źródło: opracowanie własne.}
\label{tab:przykladowe-wpisy-osi-czasu}
\scriptsize
\resizebox{\textwidth}{!}{%
\begin{tabular}{ccccrrrrrrr}
\toprule
\textbf{Klatka} &
\textbf{Czas [s]} &
\textbf{Poza} &
\textbf{Strona} &
\textbf{\texttt{knee\_angle}} &
\textbf{\texttt{hip\_angle}} &
\textbf{\texttt{shoulder\_angle}} &
\textbf{\texttt{back\_inclination}} &
\textbf{\texttt{bar\_x}} &
\textbf{\texttt{bar\_y}} &
\textbf{\texttt{bar\_visibility}} \\
\midrule
40 & 1{,}33 & tak & left & 82{,}89 & 28{,}94 & 63{,}72 & 82{,}96 & 0{,}50 & 0{,}68 & 0{,}45 \\
44 & 1{,}47 & tak & left & 108{,}28 & 42{,}40 & 58{,}69 & 79{,}97 & 0{,}50 & 0{,}66 & 0{,}42 \\
48 & 1{,}60 & tak & left & 120{,}34 & 59{,}30 & 51{,}92 & 71{,}78 & 0{,}50 & 0{,}64 & 0{,}11 \\
52 & 1{,}73 & tak & left & 145{,}80 & 90{,}70 & 41{,}47 & 59{,}75 & 0{,}50 & 0{,}60 & 0{,}16 \\
56 & 1{,}87 & tak & left & 156{,}09 & 123{,}38 & 30{,}51 & 42{,}04 & 0{,}50 & 0{,}58 & 0{,}25 \\
60 & 2{,}00 & tak & left & 174{,}38 & 169{,}59 & 20{,}59 & 15{,}76 & 0{,}49 & 0{,}57 & 0{,}14 \\
64 & 2{,}13 & tak & left & 170{,}24 & 170{,}94 & 19{,}10 & 7{,}13 & 0{,}48 & 0{,}57 & 0{,}22 \\
68 & 2{,}27 & tak & left & 177{,}92 & 161{,}81 & 15{,}46 & 1{,}78 & 0{,}47 & 0{,}56 & 0{,}28 \\
72 & 2{,}40 & tak & left & 175{,}96 & 161{,}30 & 15{,}53 & 1{,}10 & 0{,}48 & 0{,}57 & 0{,}21 \\
76 & 2{,}53 & tak & left & 179{,}24 & 160{,}83 & 16{,}00 & 2{,}28 & 0{,}47 & 0{,}57 & 0{,}20 \\
\bottomrule
\end{tabular}%
}
\end{table}

Wartości przedstawione w tabeli pokazują zmianę parametrów w trakcie przechodzenia z dolnej części ruchu do pozycji wyprostu. Widoczny jest wzrost wartości kątów \texttt{knee\_angle} oraz \texttt{hip\_angle}, co odpowiada prostowaniu kończyn dolnych i bioder w czasie podnoszenia ciężaru. Jednocześnie należy zwrócić uwagę na relatywnie niskie wartości parametru \texttt{bar\_visibility}. Nie oznaczają one jednak niezależnej pewności wykrycia sztangi, ponieważ system nie rozpoznaje gryfu jako osobnego obiektu. Wartość ta jest obliczana na podstawie widoczności nadgarstków, które w nagraniu z perspektywy bocznej mogą być częściowo zasłonięte przez ciało, sztangę lub talerze. Z tego powodu pozycja sztangi traktowana jest w systemie jako przybliżenie, a nie jako wynik osobnej detekcji obiektu.

\subsection{Wykrywanie powtórzeń i przypisywanie faz ruchu}

Po przygotowaniu osi czasu analizy możliwe jest przejście do kolejnego etapu przetwarzania, czyli wykrycia pojedynczych powtórzeń martwego ciągu. Za tę część działania systemu odpowiada moduł \texttt{repetition\_detection.py}, w którym zaimplementowano funkcje odpowiedzialne za wykrywanie powtórzeń oraz przypisywanie faz ruchu do osi czasu analizy. Na tym etapie system nie korzysta już bezpośrednio z obrazu wideo, lecz z danych liczbowych wyznaczonych podczas analizy kolejnych klatek. Podstawową informacją wykorzystywaną do określenia przebiegu powtórzenia jest pionowa pozycja sztangi, przybliżona wcześniej na podstawie położenia nadgarstków.

W pierwszej kolejności z osi czasu wybierane są wpisy z wykrytą sylwetką ćwiczącego oraz pozycją sztangi. Istotne jest pominięcie klatek bez poprawnej detekcji, ponieważ brak danych uniemożliwia określenie przebiegu ruchu z wystarczającą pewnością.

Dla każdego poprawnego wpisu zapisywany jest numer klatki, czas nagrania oraz pionowa współrzędna środka sztangi. Wartość ta jest następnie wykorzystywana do odtworzenia przebiegu ruchu w czasie. W przyjętym układzie współrzędnych większa wartość \texttt{bar\_y} oznacza niższe położenie sztangi. Wartość \texttt{0} odpowiada górnej krawędzi obrazu, a \texttt{1} dolnej.

Ponieważ pozycja sztangi jest wyznaczana pośrednio na podstawie punktów kluczowych sylwetki, jej wartości mogą podlegać niewielkim wahaniom pomiędzy kolejnymi klatkami \cite{wade2022markerless,mundt2024twodimensional}. Z tego względu przed wykrywaniem powtórzeń przebieg \texttt{bar\_y} jest wygładzany za pomocą średniej kroczącej, co ogranicza wpływ niewielkich wahań pomiędzy kolejnymi próbkami.

Następnie na wygładzonym przebiegu ruchu wyszukiwane są ekstrema lokalne. W kontekście martwego ciągu maksimum lokalne współrzędnej \texttt{bar\_y} odpowiada dolnemu położeniu sztangi, a minimum lokalne odpowiada jej położeniu górnemu. Dla każdego analizowanego punktu sprawdzane jest więc, czy jego wartość jest większa lub mniejsza od wartości znajdujących się w określonym sąsiedztwie. Wielkość tego sąsiedztwa została zdefiniowana jako stała konfiguracyjna, dzięki czemu możliwe jest dostosowanie czułości detekcji do charakteru analizowanego nagrania.

Wykrycie pojedynczego powtórzenia opiera się na sekwencji trzech kolejnych ekstremów: dolnego położenia sztangi, górnego położenia sztangi oraz ponownego dolnego położenia sztangi. Taki układ odpowiada pełnemu przebiegowi jednego powtórzenia martwego ciągu, obejmującemu fazę podnoszenia ciężaru oraz fazę jego opuszczania. Jeżeli kolejne ekstrema nie tworzą takiej sekwencji, są pomijane, a algorytm przechodzi do dalszych punktów przebiegu.

Przykładowy przebieg pionowej pozycji środka gryfu sztangi w czasie przedstawiono na rysunku~\ref{fig:wykres-zmiana-gryfu}.

\begin{figure}[H]
\centering
\includegraphics[width=\textwidth]{figures/wykres_zmiana_gryfu.png}
\caption{Przebieg pionowej pozycji środka gryfu sztangi w czasie. Źródło: opracowanie własne.}
\label{fig:wykres-zmiana-gryfu}
\end{figure}

Na rysunku~\ref{fig:wykres-zmiana-gryfu} widoczny jest cykliczny przebieg wartości \texttt{bar\_y}, odpowiadający kolejnym powtórzeniom martwego ciągu. Spadek wartości tej współrzędnej oznacza ruch gryfu w górę, natomiast jej wzrost odpowiada fazie opuszczania ciężaru. Lokalne maksima i minima przebiegu umożliwiają wyznaczenie charakterystycznych punktów powtórzenia, czyli dolnej pozycji, górnej pozycji oraz powrotu do pozycji dolnej. W tym rozdziale wykres służy głównie do przedstawienia sposobu działania implementacji, natomiast szczegółowa ocena poprawności detekcji została omówiona w części poświęconej testom systemu.

Wystąpienie sekwencji dolnego, górnego i dolnego położenia sztangi nie świadczy jeszcze o wykryciu pełnego powtórzenia. Dodatkowo sprawdzany jest pionowy zakres ruchu sztangi oraz czas trwania ruchu. Zakres musi przekroczyć minimalną zadaną wartość zarówno podczas podnoszenia, jak i opuszczania. Warunek ten ogranicza możliwość zakwalifikowania niewielkich wahań położenia jako osobnego powtórzenia. Dodatkowo sprawdzany jest minimalny czas trwania ruchu, co pozwala odrzucić zbyt krótkie fragmenty.

Ważnym elementem implementacji jest osobne wyznaczenie początku fazy podnoszenia. Pierwsze dolne ekstremum wskazuje moment, w którym sztanga znajduje się nisko, jednak nie zawsze oznacza dokładny moment rozpoczęcia ruchu w górę. Z tego względu algorytm wyszukuje pierwszą klatkę, w której sztanga przesunęła się w górę o określoną minimalną wartość względem pozycji dolnej. Dopiero ten punkt traktowany jest jako właściwy początek fazy podnoszenia.

Przed fazą podnoszenia wyznaczana jest krótka faza przygotowawcza. Obejmuje ona fragment poprzedzający rozpoczęcie podnoszenia, w którym sztanga znajduje się w wyznaczonej dolnej strefie ruchu. Analizowany fragment jest ograniczony do maksymalnie jednej sekundy przed początkiem fazy podnoszenia.

Dla każdego wykrytego przez system powtórzenia zapisywane są najważniejsze punkty w czasie oraz odpowiadające im numery klatek. Punkty te to kolejno: początek fazy przygotowania, początek fazy podnoszenia, górna pozycja sztangi oraz koniec fazy opuszczania. Na podstawie różnic w wartościach przypisanych do tych punktów obliczane są czasy poszczególnych faz ruchu. Dodatkowo zapisywany jest pionowy zakres ruchu sztangi osobno dla fazy podnoszenia i opuszczania.

Po wykryciu powtórzeń następuje przypisanie ich do osi czasu analizy. W implementacji odpowiada za to funkcja \texttt{assign\_repetitions\_to\_timeline}, która uzupełnia wpisy osi czasu o numer powtórzenia oraz nazwę aktualnej fazy ruchu. Każdy wpis osi czasu otrzymuje domyślnie informację, że nie należy do żadnego powtórzenia oraz znajduje się w stanie bezczynności. Następnie dla każdego wykrytego powtórzenia sprawdzany jest numer klatki i na tej podstawie wpisowi przypisywany jest numer powtórzenia oraz jedna z faz ruchu. Klatki od początku przygotowania do początku podnoszenia oznaczane są jako faza przygotowania, klatki od początku podnoszenia do górnej pozycji sztangi jako faza podnoszenia, natomiast klatki od górnej pozycji sztangi do końca opuszczania jako faza opuszczania.

W ten sposób oś czasu zostaje rozszerzona o informację dotyczącą przynależności każdej klatki do konkretnego powtórzenia. Jest to istotne dla kolejnego etapu działania systemu, w którym ocena techniki jest przeprowadzana na podstawie danych opisujących całe powtórzenie, a nie pojedyncze klatki analizowanego materiału wideo.

\subsection{Tworzenie podsumowań pojedynczych powtórzeń}

Po wykryciu powtórzeń oraz przypisaniu faz ruchu do osi czasu możliwe jest przygotowanie zbiorczego opisu każdego powtórzenia. Za ten etap działania systemu odpowiada moduł \texttt{repetition\_summary.py}. Jego zadaniem jest przekształcenie danych zapisanych dla wielu klatek w jeden zestaw parametrów opisujących pojedyncze powtórzenie martwego ciągu.

Pierwszym etapem tworzenia podsumowania jest wybranie z osi czasu tych wpisów, które należą do analizowanego powtórzenia. Uwzględniane są wyłącznie klatki z wykrytą sylwetką oraz dostępną strukturą metryk biomechanicznych. Poszczególne wartości równe \texttt{None} są pomijane podczas tworzenia zestawień dla konkretnych metryk. Następnie wybrane wpisy są dzielone zgodnie z przypisaną fazą ruchu na fragment przygotowania, podnoszenia oraz opuszczania.

Podział na fazy umożliwia obliczenie czasu trwania poszczególnych części ruchu. Dla każdej fazy wyznaczany jest czas między pierwszym i ostatnim wpisem należącym do tej fazy. W ten sposób system określa czas przygotowania, czas podnoszenia oraz czas opuszczania. Dodatkowo obliczany jest całkowity czas trwania powtórzenia, liczony od początku fazy podnoszenia do końca fazy opuszczania. Takie podejście pozwala pominąć fragment przygotowania przy wyznaczaniu właściwego czasu ruchu roboczego.

Następnie dla danego powtórzenia analizowane są wartości wybranych metryk biomechanicznych. System pobiera z kolejnych klatek wartości kąta biodra, kąta kolana oraz kąta pochylenia tułowia. Dla każdej z tych metryk wyznaczana jest wartość minimalna, maksymalna oraz zakres zmiany w trakcie całego powtórzenia. Analogicznie analizowana jest pionowa pozycja sztangi, dla której obliczany jest minimalny i maksymalny poziom położenia oraz zakres ruchu. Zakresy zmian tych parametrów są następnie wykorzystywane podczas oceny techniki oraz porównywania kolejnych powtórzeń.

Oprócz wartości minimalnych i maksymalnych system zapisuje również wartości metryk w wybranych momentach powtórzenia. Szczególne znaczenie mają początek fazy podnoszenia oraz górna pozycja sztangi. Dla początku podnoszenia zapisywane są między innymi wartości kąta biodra, kąta kolana oraz kąta pochylenia tułowia. Dla pozycji górnej zapisywane są analogiczne wartości, które mogą być później wykorzystane do oceny stopnia wyprostu biodra i kolan.

W module podsumowania uwzględniono również informacje o położeniu wybranych punktów kluczowych w momencie rozpoczęcia podnoszenia. System odczytuje położenie barku oraz biodra po stronie sylwetki wybranej wcześniej do analizy. Na tej podstawie określa, czy bark znajduje się niżej niż biodro w układzie współrzędnych obrazu. Informacja ta może być przydatna przy późniejszej ocenie ustawienia tułowia na początku ruchu.

Wynikiem działania funkcji \texttt{summarize\_repetition} jest słownik zawierający zestaw parametrów opisujących jedno powtórzenie. Obejmuje on między innymi numery i czasy charakterystycznych punktów ruchu, czasy trwania faz, zakres pionowego ruchu sztangi, zmiany kątów oraz wartości wybranych metryk w pozycji początkowej i końcowej. Następnie funkcja \texttt{summarize\_repetitions} wykonuje ten sam proces dla wszystkich wykrytych powtórzeń i zwraca listę ich podsumowań.

Wybrane parametry tworzone w ramach podsumowania pojedynczego powtórzenia przedstawiono w tabeli~\ref{tab:parametry-podsumowania-powtorzenia}.

\begin{table}[H]
\centering
\caption{Wybrane parametry podsumowania pojedynczego powtórzenia. Źródło: opracowanie własne.}
\label{tab:parametry-podsumowania-powtorzenia}
\vspace{0.3cm}
\small
\begin{tabular}{p{5.6cm}p{8.4cm}}
\toprule
\addlinespace[0.3em]
\textbf{Parametr} & \textbf{Znaczenie} \\
\midrule
\texttt{rep\_number} & Numer wykrytego powtórzenia w analizowanym nagraniu. \\

\texttt{lifting\_duration\_seconds} & Czas trwania fazy podnoszenia ciężaru, liczony od rozpoczęcia ruchu sztangi w górę do osiągnięcia pozycji górnej. \\

\texttt{lowering\_duration\_seconds} & Czas trwania fazy opuszczania ciężaru, liczony od pozycji górnej do końca opuszczania sztangi. \\

\texttt{duration\_seconds} & Całkowity czas trwania właściwego ruchu w ramach powtórzenia, bez uwzględniania fazy przygotowania. \\

\texttt{bar\_y\_range} & Zakres pionowego ruchu środka gryfu sztangi w analizowanym powtórzeniu. \\

\texttt{hip\_angle\_change} & Różnica pomiędzy maksymalną i minimalną wartością kąta biodra w trakcie powtórzenia. \\

\texttt{knee\_angle\_change} & Różnica pomiędzy maksymalną i minimalną wartością kąta kolana w trakcie powtórzenia. \\

\texttt{back\_angle\_change} & Różnica pomiędzy maksymalną i minimalną wartością kąta pochylenia tułowia w trakcie powtórzenia. \\

\texttt{top\_hip\_angle} & Wartość kąta biodra w górnej pozycji sztangi. \\

\texttt{top\_knee\_angle} & Wartość kąta kolana w górnej pozycji sztangi. \\

\texttt{lifting\_start\_hip\_angle} & Wartość kąta biodra w momencie rozpoczęcia fazy podnoszenia. \\

\texttt{lifting\_start\_knee\_angle} & Wartość kąta kolana w momencie rozpoczęcia fazy podnoszenia. \\

\texttt{shoulder\_below\_hip\_at\_start} & Informacja określająca, czy w momencie rozpoczęcia podnoszenia bark znajduje się niżej niż biodro w układzie współrzędnych obrazu. \\
\bottomrule
\end{tabular}
\end{table}

Tak przygotowane podsumowania stanowią warstwę pośrednią pomiędzy analizą klatka po klatce a oceną techniki. Dzięki nim kolejne etapy systemu nie muszą ponownie przetwarzać całej osi czasu nagrania. Zamiast tego mogą korzystać z uporządkowanego zestawu parametrów opisujących każde powtórzenie w jednolity sposób.

\subsection{Ocena techniki wykonania powtórzeń}

Po przygotowaniu podsumowań pojedynczych powtórzeń system przechodzi do etapu oceny techniki. Za tę część odpowiada moduł \texttt{technique\_evaluation.py}. Nie analizuje on już całej osi czasu nagrania, lecz korzysta z wcześniej obliczonych parametrów opisujących każde powtórzenie. Dzięki temu ocena wykonywana jest na uporządkowanych danych, a nie bezpośrednio na pojedynczych klatkach.

Ocena techniki została oparta na zestawie reguł decyzyjnych. Każda z nich sprawdza wybrany element ruchu, co pozwala wskazać, które cechy wykonania wpłynęły na uzyskany wynik.

Pierwszym elementem poddawanym ocenie jest pozycja końcowa ruchu. W martwym ciągu górna pozycja oznacza wyprost bioder oraz kolan. Z tego powodu system bierze pod uwagę wartości zmiennych \texttt{top\_hip\_angle} oraz \texttt{top\_knee\_angle}. Jeżeli wartości te są niższe od przyjętych progów, system generuje ostrzeżenie dotyczące niepełnego wyprostu biodra lub kolana. System nie ogranicza się więc do samego wykrycia ruchu gryfu w górę, ale uwzględnia również pozycję końcową osoby ćwiczącej.

Kolejnym elementem jest zakres pionowego ruchu gryfu. Parametr \texttt{bar\_y\_range} określa różnicę pomiędzy dolnym i górnym położeniem środka gryfu w danym powtórzeniu. Przed wykonaniem oceny dla całej serii wyznaczana jest mediana tego parametru. Zakres danego powtórzenia jest następnie porównywany z wartością odpowiadającą określonej części mediany zakresu ruchu serii. Pozwala to wykryć powtórzenia, w których zakres ruchu był wyraźnie mniejszy niż w pozostałych.

Przed właściwą oceną system uzupełnia podsumowania o kontekst całej serii. Wykorzystywana jest mediana zakresu ruchu gryfu oraz mediana kąta kolana w pozycji startowej. Takie podejście ogranicza wpływ pojedynczych nietypowych powtórzeń na wartości odniesienia oraz pozwala porównywać kolejne wykonania w obrębie jednego nagrania.

Osobno analizowana jest pozycja początkowa powtórzenia. Jej charakterystyka określana jest na podstawie kąta kolana w odniesieniu do mediany tej wartości dla całej serii. Pozwala to wskazać, czy dane powtórzenie rozpoczynało się z pozycji niższej, wyższej lub zbliżonej do pozostałych. W komunikacie pomocniczo uwzględniane są również wartości kąta biodra oraz pochylenia tułowia. Informacja ta ma charakter opisowy i nie wpływa bezpośrednio na klasyfikację poprawności powtórzenia.

Moduł oceny bierze również pod uwagę relację położenia barku i biodra w momencie rozpoczęcia ruchu. Jeżeli bark znajduje się niżej niż biodro, system zwraca ostrzeżenie dotyczące ustawienia tułowia. Warunek ten nie stanowi bezpośredniej oceny krzywizny kręgosłupa, lecz służy do wskazania fragmentu wykonania wymagającego dodatkowej uwagi.

Dodatkowo generowane są informacje dotyczące zmiany pochylenia tułowia oraz tempa wykonania powtórzenia. Zmiana pochylenia tułowia wynika z różnicy wartości kąta pleców w trakcie ruchu. Tempo jest natomiast opisane przez czas fazy podnoszenia i czas fazy opuszczania. Parametry te mają charakter informacyjny i pomagają opisać przebieg powtórzenia, ale nie wpływają bezpośrednio na jego klasyfikację jako poprawnego, częściowego lub niepoprawnego.

Każda reguła zwraca odpowiedni komunikat oraz status. Dla poszczególnych sprawdzeń wykorzystywane są trzy statusy: pozytywny, informacyjny oraz ostrzegawczy. Status pozytywny wskazuje, że w danym elemencie nie wykryto istotnej rozbieżności. Status informacyjny oznacza, że dana wartość ma charakter pomocniczy i nie stanowi bezpośredniej oceny poprawności wykonania ćwiczenia. Status ostrzegawczy wskazuje natomiast element, który wymaga dodatkowej uwagi.

Niezależnie od statusów poszczególnych sprawdzeń system określa również klasyfikację całego powtórzenia. Przy jej wyznaczaniu brane są pod uwagę trzy podstawowe elementy: wyprost biodra, wyprost kolana oraz zakres pionowego ruchu gryfu. Pozostałe reguły mają charakter pomocniczy lub ostrzegawczy, ale nie wpływają bezpośrednio na tę klasyfikację.

Wybrane reguły wykorzystywane podczas oceny techniki pojedynczego powtórzenia przedstawiono w tabeli~\ref{tab:reguly-oceny-techniki}.

\begin{table}[H]
\centering
\caption{Wybrane reguły wykorzystywane podczas oceny techniki powtórzenia. Źródło: opracowanie własne.}
\label{tab:reguly-oceny-techniki}
\vspace{0.3cm}
\footnotesize
\renewcommand{\arraystretch}{1.25}
\begin{tabular}{|p{3.8cm}|p{2.8cm}|p{7.0cm}|}
\hline
\textbf{Element oceny} & \textbf{Typ reguły} & \textbf{Znaczenie} \\
\hline

Wyprost biodra w pozycji końcowej & Oceniająca & Reguła sprawdza wartość parametru \texttt{top\_hip\_angle}. Wartość poniżej przyjętego progu powoduje wygenerowanie ostrzeżenia dotyczącego niepełnego wyprostu biodra. \\
\hline

Wyprost kolan w pozycji końcowej & Oceniająca & Reguła sprawdza wartość parametru \texttt{top\_knee\_angle}. Wartość poniżej przyjętego progu powoduje wygenerowanie ostrzeżenia dotyczącego niepełnego wyprostu kolana. \\
\hline

Zakres pionowego ruchu gryfu & Oceniająca & Reguła wykorzystuje parametr \texttt{bar\_y\_range} i porównuje go z wartością odniesienia wyznaczoną na podstawie mediany zakresu ruchu w całej serii. \\
\hline

Pozycja startowa powtórzenia & Informacyjna & Reguła porównuje kąt kolana w momencie rozpoczęcia podnoszenia z medianą tej wartości dla serii. Pomocniczo prezentowane są również wartości kąta biodra i pochylenia tułowia. \\
\hline

Ustawienie barku względem biodra & Ostrzegawcza & Reguła wykorzystuje informację zapisaną w parametrze \texttt{shoulder\_below\_hip\_at\_start}. Może wskazywać na niekorzystne ustawienie tułowia na początku podnoszenia. \\
\hline

Zmiana pochylenia tułowia & Informacyjna & Reguła opisuje zmianę kąta pochylenia tułowia w trakcie powtórzenia. Informacja ta pomaga określić, jak zmieniało się ustawienie pleców podczas ruchu. \\
\hline

Tempo wykonania powtórzenia & Informacyjna & Reguła wykorzystuje czas fazy podnoszenia i opuszczania. Parametr ten pozwala lepiej opisać przebieg pojedynczego powtórzenia. \\
\hline

\end{tabular}
\end{table}

Na podstawie trzech głównych reguł oceniana jest klasyfikacja powtórzenia. Jeżeli żadna z nich nie wskazuje problemu, powtórzenie zostaje zaklasyfikowane jako poprawne. W przypadku wykrycia jednego lub dwóch problemów otrzymuje klasyfikację częściową. Jeżeli problemy zostaną wykryte we wszystkich trzech ocenianych elementach, powtórzenie zostaje zaklasyfikowane jako niepoprawne.

Główna funkcja odpowiedzialna za dokonanie oceny pojedynczego powtórzenia nosi nazwę \texttt{evaluate\_repetition}. Przyjmuje ona podsumowanie jednego powtórzenia i zwraca zestaw komunikatów, ogólny status, klasyfikację powtórzenia oraz powody wpływające na wynik. Funkcja \texttt{evaluate\_repetitions} wykonuje ten proces dla wszystkich powtórzeń w serii. Przed rozpoczęciem oceny uzupełnia ona podsumowania o kontekst serii, a następnie przypisuje wynik do każdego wykrytego powtórzenia.

\subsection{Ocena spójności serii}

Po ocenie pojedynczych powtórzeń system wykonuje jeszcze jeden etap analizy, którego celem jest sprawdzenie spójności całej serii. Za tę część odpowiedzialna jest funkcja o nazwie \texttt{evaluate\_repetition\_consistency} umieszczona w module \texttt{technique\_evaluation.py}. W przeciwieństwie do poprzedniego etapu nie analizuje ona osobno każdego powtórzenia, lecz porównuje powtórzenia między sobą.

Do wyboru podsumowań wykorzystywanych podczas tego porównania służy funkcja pomocnicza \texttt{get\_summaries\_for\_consistency}. W pierwszej kolejności brane są pod uwagę powtórzenia oznaczone wcześniej jako poprawne. Jeżeli dostępne są co najmniej dwa takie powtórzenia, to właśnie one są wykorzystywane do oceny spójności. W przeciwnym przypadku system korzysta ze wszystkich dostępnych podsumowań.

Ocena spójności serii ma charakter pomocniczy. Nie służy do wykrywania pojedynczego błędu technicznego, lecz do określenia, jak bardzo poszczególne powtórzenia różnią się od siebie. Niewielkie różnice pomiędzy kolejnymi powtórzeniami są naturalne. Celem tej części analizy jest wskazanie sytuacji, w których zmienność wybranych parametrów przekracza przyjęte wartości progowe.

Pierwszym porównywanym przez program parametrem jest zakres ruchu gryfu sztangi. System oblicza różnicę pomiędzy maksymalną i minimalną wartością parametru \texttt{bar\_y\_range} występującą w analizowanych powtórzeniach. Na podstawie tej różnicy można określić, czy kolejne powtórzenia miały podobny zakres ruchu, czy też różniły się od siebie w stopniu przekraczającym przyjęty próg. Duża rozbieżność może wskazywać na brak powtarzalności zakresu ruchu w obrębie serii.

Drugim elementem jest czas trwania powtórzeń. W tym przypadku system porównuje wartości parametru \texttt{duration\_seconds}. Różnice czasu wykonania nie muszą oznaczać błędu technicznego, dlatego ten element ma charakter informacyjny. Może jednak pokazać, że tempo wykonania kolejnych powtórzeń nie było jednakowe.

Ostatnim porównaniem jest sprawdzenie wartości kątów w pozycji końcowej. System oblicza różnice dotyczące parametrów \texttt{top\_hip\_angle} oraz \texttt{top\_knee\_angle}. Na podstawie tych wartości możliwe jest określenie, czy ćwiczący kończył kolejne powtórzenia w podobnej pozycji. Duże różnice wskazują na brak powtarzalności pozycji końcowej pomiędzy powtórzeniami.

Parametry wykorzystywane podczas oceny spójności serii przedstawiono w tabeli~\ref{tab:parametry-spojnosci-serii}.

\begin{table}[H]
\centering
\caption{Parametry wykorzystywane podczas oceny spójności serii. Źródło: opracowanie własne.}
\label{tab:parametry-spojnosci-serii}
\vspace{0.3cm}
\footnotesize
\renewcommand{\arraystretch}{1.25}
\begin{tabular}{|p{3.6cm}|p{6.3cm}|p{3.7cm}|}
\hline
\textbf{Parametr} & \textbf{Porównywana wartość} & \textbf{Charakter wyniku} \\
\hline

\texttt{bar\_y\_range} & Różnica zakresu pionowego ruchu gryfu pomiędzy powtórzeniami w serii. & Oceniający \\
\hline

\texttt{duration\_seconds} & Różnica czasu trwania powtórzeń, liczona na podstawie czasu ruchu roboczego. & Informacyjny \\
\hline

\texttt{top\_hip\_angle} & Różnica wartości kąta biodra w górnej pozycji powtórzeń. & Ostrzegawczy \\
\hline

\texttt{top\_knee\_angle} & Różnica wartości kąta kolana w górnej pozycji powtórzeń. & Ostrzegawczy \\
\hline

\end{tabular}
\end{table}

Dla każdego sprawdzanego elementu tworzony jest osobny komunikat. System oblicza różnicę pomiędzy największą i najmniejszą wartością danego parametru, a następnie porównuje ją z odpowiednim progiem zapisanym w konfiguracji. Jeżeli wartość mieści się w dopuszczalnym zakresie, sprawdzenie otrzymuje status pozytywny. W przypadku przekroczenia progu dla zakresu ruchu gryfu lub kątów w pozycji końcowej zwracany jest status ostrzegawczy. Przekroczenie progu dotyczącego czasu trwania powtórzeń powoduje natomiast zwrócenie statusu informacyjnego.

Na wynik oceny spójności całej serii składają się status ogólny oraz lista poszczególnych sprawdzeń. Jeżeli żaden z elementów nie otrzyma statusu ostrzegawczego, całej ocenie przypisywany jest status pozytywny. Jeżeli co najmniej jedno sprawdzenie dotyczące zakresu ruchu gryfu lub pozycji końcowej otrzyma status ostrzegawczy, cała ocena zostaje oznaczona jako wymagająca uwagi. Sama różnica czasu trwania powtórzeń nie wpływa na końcowy status ostrzegawczy serii.

Ten etap uzupełnia ocenę pojedynczych powtórzeń o szerszy kontekst. Dzięki niemu użytkownik może zobaczyć nie tylko, czy dane powtórzenie zostało wykonane poprawnie, ale również czy cała seria była wykonywana w powtarzalny sposób. Jest to szczególnie przydatne przy analizie kilku powtórzeń w jednym nagraniu, ponieważ pozwala zauważyć różnice, które nie zawsze są widoczne przy ocenie każdego powtórzenia osobno.

\subsection{Integracja modułu analizy z aplikacją serwerową}

Po zakończeniu opisu modułu analizy nagrania wideo należy przedstawić sposób jego połączenia z aplikacją serwerową. W projekcie moduł analizy nie został wydzielony jako osobny mikroserwis. Działa on w tym samym środowisku uruchomieniowym co aplikacja serwerowa. Takie rozwiązanie ogranicza liczbę elementów infrastruktury potrzebnych do uruchomienia systemu i jest wystarczające dla zakresu pracy inżynierskiej.

Za integrację części serwerowej z modułem analizy odpowiada \texttt{analysis\_service.py}, znajdujący się w aplikacji backendowej. Pełni on rolę warstwy pośredniej pomiędzy kodem odpowiedzialnym za obsługę żądania użytkownika a właściwym modułem analizy wideo. Dzięki temu widoki aplikacji serwerowej nie muszą bezpośrednio znać szczegółów działania klasy \texttt{DeadliftAnalyzer}. Ich zadaniem jest obsługa przesłanego pliku oraz przekazanie jego ścieżki do odpowiedniej funkcji serwisowej.

W pliku \texttt{analysis\_service.py} wyznaczana jest ścieżka do głównego katalogu projektu oraz do katalogu \texttt{ml\_model}. Następnie katalog ten jest dodawany do listy ścieżek importu Pythona. Pozwala to aplikacji serwerowej korzystać z modułów znajdujących się poza bezpośrednią strukturą aplikacji Django. W ten sposób backend może zaimportować klasę \texttt{DeadliftAnalyzer} bez konieczności kopiowania kodu analizy do katalogu aplikacji serwerowej.

Główna funkcja wykorzystywana w tej warstwie to \texttt{run\_deadlift\_analysis}. Przyjmuje ona ścieżkę do pliku wideo zapisanego po stronie serwera. Następnie tworzy instancję klasy \texttt{DeadliftAnalyzer} i wywołuje analizę dla otrzymanej ścieżki. Oznacza to, że aplikacja serwerowa nie przekazuje do modułu analizy samego pliku wideo w postaci binarnej, lecz lokalizację pliku już zapisanego na dysku.

Przepływ wywołania analizy nagrania z poziomu aplikacji serwerowej przedstawiono na rysunku~\ref{fig:diagram-integracji-modulu-analizy}.

\begin{figure}[H]
\centering
\includegraphics[width=\textwidth]{figures/diagram_integracji_modulu_analizy.png}
\caption{Przepływ wywołania analizy nagrania z poziomu aplikacji serwerowej. Źródło: opracowanie własne.}
\label{fig:diagram-integracji-modulu-analizy}
\end{figure}

Diagram przedstawia bezpośrednie wywołanie modułu analitycznego przez warstwę serwisową aplikacji backendowej oraz przekazanie wyniku do dalszego zapisu. Moduł analizy nie musi dzięki temu obsługiwać żądań HTTP ani komunikacji sieciowej.

Wynik zwrócony przez klasę \texttt{DeadliftAnalyzer} jest przygotowywany do zapisu w polu \texttt{JSONField} modelu analizy. W tym celu w pliku \texttt{analysis\_service.py} zaimplementowano funkcję \texttt{make\_json\_serializable}. Jej zadaniem jest przejście przez zwróconą strukturę danych i przekształcenie wartości, które nie mogą zostać bezpośrednio zapisane w formacie JSON.

Funkcja \texttt{make\_json\_serializable} obsługuje między innymi słowniki, listy i krotki oraz wartości pochodzące z bibliotek wykorzystywanych w module analizy. Krotki zamieniane są na listy, natomiast wartości posiadające metodę \texttt{item} przekształcane są do standardowych typów języka Python. Jest to istotne, ponieważ część obliczeń wykonywanych podczas analizy może zwracać typy danych, które nie mogą zostać bezpośrednio zapisane jako dane JSON.

Przyjęta struktura kodu nie wyklucza późniejszego wydzielenia modułu analizy do osobnej usługi. W obecnej implementacji nie było to jednak konieczne.

Dzięki zastosowaniu warstwy serwisowej aplikacja serwerowa pozostaje oddzielona od szczegółów działania algorytmów analizy. Widok odpowiedzialny za obsługę przesłanego nagrania może wywołać funkcję integracyjną i otrzymać gotowy wynik analizy bez bezpośredniego korzystania z wewnętrznych elementów modułu analitycznego.

\subsection{Implementacja aplikacji serwerowej}

Aplikacja serwerowa, określana dalej również jako backend, odpowiada za obsługę żądań wysyłanych z części klienckiej, zapis przesłanego nagrania, uruchomienie analizy oraz zwrócenie jej wyniku. W projekcie została ona zaimplementowana z wykorzystaniem Django, czyli szkieletu aplikacyjnego przeznaczonego do tworzenia aplikacji internetowych w języku Python.

Najważniejszym obiektem wchodzącym w skład backendu jest model \texttt{Analysis}. Reprezentuje on pojedyncze zlecenie analizy wywołane przez aplikację kliencką. W rekordzie przechowywany jest plik wideo, aktualny status analizy, wynik zwrócony przez system oraz ewentualny komunikat błędu. Dodatkowo zapisywane są znaczniki czasu odpowiadające utworzeniu rekordu, rozpoczęciu analizy oraz jej zakończeniu.

Nagranie przechowywane jest za pomocą pola \texttt{FileField}. Plik zapisywany jest w katalogu mediów aplikacji, natomiast w rekordzie modelu przechowywane jest odwołanie pozwalające na dostęp do jego lokalizacji. Ścieżka do zapisanego pliku jest następnie wykorzystywana podczas uruchamiania modułu analizy.

Model \texttt{Analysis} wykorzystuje kilka statusów opisujących aktualny stan analizy. Status \texttt{PENDING} oznacza, że rekord został utworzony, ale analiza nie została jeszcze rozpoczęta. Status \texttt{PROCESSING} informuje, że system przetwarza przesłany materiał. Po poprawnym zakończeniu analizy ustawiany jest status \texttt{COMPLETED}. Jeżeli podczas przetwarzania wystąpi błąd, rekord otrzymuje status \texttt{FAILED}, a szczegóły problemu zapisywane są w polu \texttt{error\_message}.

Przepływ zmian statusu rekordu analizy przedstawiono na rysunku~\ref{fig:diagram-stanow-analizy}.

\begin{figure}[H]
\centering
\includegraphics[width=0.85\textwidth]{figures/diagram_stanow_rekordow.png}
\caption{Diagram stanów rekordu analizy w aplikacji serwerowej. Źródło: opracowanie własne.}
\label{fig:diagram-stanow-analizy}
\end{figure}

Wynik analizy przechowywany jest w polu \texttt{result} typu \texttt{JSONField}. Pozwala ono zapisać złożoną strukturę danych zwracaną przez moduł analizy bez konieczności tworzenia osobnych pól modelu dla każdego elementu wyniku.

Za serializację danych modelu \texttt{Analysis} odpowiada klasa \texttt{AnalysisSerializer}. Udostępnia ona podstawowe pola modelu oraz dodatkowe informacje pomocnicze, takie jak nazwa przesłanego pliku i czas trwania analizy. Część pól została oznaczona jako tylko do odczytu, ponieważ nie powinna być ustawiana bezpośrednio przez użytkownika. Dotyczy to między innymi statusu, wyniku analizy, komunikatu błędu oraz znaczników czasu.

Głównym widokiem odpowiedzialnym za rozpoczęcie analizy jest \texttt{create\_analysis}. Przyjmuje ona żądanie zawierające plik wideo przesłany przez użytkownika. W pierwszej kolejności dane wejściowe są przekazywane do serializera. Jeżeli przejdą jego walidację, tworzony jest nowy rekord modelu \texttt{Analysis} ze statusem \texttt{PENDING}. Następnie status zostaje zmieniony na \texttt{PROCESSING}, a w polu \texttt{started\_at} zapisywany jest czas rozpoczęcia przetwarzania.

Po przygotowaniu rekordu widok wywołuje funkcję \texttt{run\_deadlift\_analysis}, przekazując do niej ścieżkę do zapisanego pliku wideo. Po zakończeniu analizy zwrócony wynik zostaje przypisany do pola \texttt{result}, status zmieniany jest na \texttt{COMPLETED}, a czas zakończenia zapisywany jest w polu \texttt{completed\_at}.

Widok \texttt{create\_analysis} obsługuje również błędy powstałe podczas wykonywania analizy. W przypadku wystąpienia wyjątku status rekordu zmieniany jest na \texttt{FAILED}, treść błędu zapisywana jest w polu \texttt{error\_message}, a w polu \texttt{completed\_at} ustawiany jest czas zakończenia przetwarzania. Dzięki temu informacja o niepowodzeniu pozostaje zapisana i może zostać przekazana aplikacji klienckiej.

Drugim istotnym widokiem jest funkcja \texttt{get\_analysis}. Jej zadaniem jest pobranie istniejącego rekordu analizy na podstawie identyfikatora. Jeżeli rekord zostanie znaleziony, widok zwraca jego dane w formacie JSON. W przypadku braku rekordu o podanym identyfikatorze zwracana jest odpowiedź informująca o nieznalezieniu zasobu.

W obecnej implementacji analiza uruchamiana jest synchronicznie w ramach obsługi żądania HTTP. Oznacza to, że po przesłaniu nagrania backend wykonuje analizę i dopiero po jej zakończeniu zwraca odpowiedź do aplikacji klienckiej. Takie rozwiązanie jest wystarczające dla zakresu projektu i nie wymaga zastosowania dodatkowej kolejki zadań ani osobnego procesu roboczego.

\subsection{Implementacja aplikacji klienckiej}

Aplikacja kliencka, określana dalej również jako frontend, odpowiada za część systemu widoczną dla użytkownika. Jej głównym zadaniem jest umożliwienie przesłania nagrania wideo, uruchomienie analizy oraz przedstawienie użytkownikowi informacji zwróconych przez aplikację serwerową. W projekcie frontend został zaimplementowany z wykorzystaniem biblioteki React, która pozwala budować interfejs użytkownika z mniejszych, niezależnych fragmentów nazywanych komponentami.

Interfejs został podzielony na komponenty odpowiadające za główne elementy aplikacji, między innymi wybór pliku, instrukcję przygotowania nagrania oraz prezentację wyników. Taki podział pozwala rozdzielić odpowiedzialność poszczególnych części interfejsu i ułatwia organizację kodu aplikacji klienckiej.

Strukturę najważniejszych elementów aplikacji klienckiej przedstawiono na rysunku~\ref{fig:diagram-struktury-frontendu}.

\begin{figure}[H]
\centering
\includegraphics[width=0.85\textwidth]{figures/diagram_komponentow.png}
\caption{Struktura głównych elementów aplikacji klienckiej. Źródło: opracowanie własne.}
\label{fig:diagram-struktury-frontendu}
\end{figure}

Na rysunku~\ref{fig:diagram-struktury-frontendu} pokazano, że głównym elementem aplikacji klienckiej jest komponent \texttt{App.jsx}, który przechowuje stan aplikacji i przekazuje dane do pozostałych komponentów. Komponent \texttt{FileUploadCard} odpowiada za wybór i podstawową walidację pliku, natomiast komunikacja z aplikacją serwerową została wydzielona do pliku \texttt{analysisApi.js}. Wartości wspólne dla kilku elementów interfejsu zapisano w pliku \texttt{constants.js}.

Głównym plikiem aplikacji klienckiej jest \texttt{App.jsx}. Przechowuje on podstawowy stan aplikacji za pomocą mechanizmu \texttt{useState}. Obejmuje on wybrany plik wideo, wynik analizy, informację o trwającym przetwarzaniu oraz ewentualny komunikat błędu. Zmiana tych wartości powoduje odpowiednią aktualizację interfejsu, dzięki czemu aplikacja może reagować na wybór pliku, rozpoczęcie analizy, otrzymanie wyniku lub wystąpienie błędu.

Obsługa wyboru pliku została wydzielona do komponentu \texttt{FileUploadCard}. Komponent ten odpowiada za wyświetlenie formularza przesyłania nagrania oraz za podstawową walidację pliku po stronie przeglądarki. W pierwszej kolejności sprawdzane jest, czy użytkownik rzeczywiście wskazał plik. Następnie weryfikowane jest, czy przeglądarka rozpoznaje jego typ jako materiał wideo. Dodatkowo sprawdzany jest maksymalny rozmiar pliku określony w stałych konfiguracyjnych aplikacji.

W przypadku wybrania nieprawidłowego pliku komponent przekazuje do głównego widoku odpowiedni komunikat błędu. Dzięki temu użytkownik otrzymuje informację jeszcze przed wysłaniem żądania do backendu.

W pliku \texttt{constants.js} umieszczono wspólne wartości konfiguracyjne oraz część tekstów wykorzystywanych w interfejsie. Zdefiniowano tam między innymi adres API, maksymalny rozmiar przesyłanego pliku, ustawienia pola wyboru pliku, etykiety statusów oraz komunikaty błędów. Takie rozwiązanie ogranicza wielokrotne definiowanie tych samych wartości i pozwala wprowadzać część zmian w jednym miejscu.

Za obsługę komunikacji strony klienckiej z aplikacją serwerową odpowiada funkcja  o nazwie \texttt{uploadVideoForAnalysis}, znajdująca się w pliku \texttt{analysisApi.js}. Funkcja rozpoczyna działanie od utworzenia obiektu \texttt{FormData}, do którego dodawany jest wybrany plik wideo pod nazwą \texttt{video}. Następnie wysyłane jest żądanie typu \texttt{POST} do odpowiedniego punktu końcowego interfejsu API. Po otrzymaniu odpowiedzi funkcja odczytuje zwrócone dane i przekazuje je do głównego komponentu aplikacji. W przypadku odpowiedzi informującej o błędzie zgłaszany jest wyjątek, który może zostać obsłużony przez interfejs użytkownika.

W komponencie \texttt{App.jsx} wysłanie formularza obsługuje funkcja \texttt{handleSubmit}. Przed rozpoczęciem komunikacji z backendem funkcja sprawdza, czy użytkownik wybrał plik. Jeżeli plik nie został wskazany, w interfejsie ustawiany jest komunikat błędu. W przeciwnym przypadku aplikacja ustawia stan ładowania, usuwa wcześniejszy komunikat błędu oraz czyści poprzedni wynik analizy. Następnie wywoływana jest funkcja \texttt{uploadVideoForAnalysis}.

Po otrzymaniu poprawnej odpowiedzi z backendu wynik analizy zostaje zapisany w stanie aplikacji. Jeżeli podczas wysyłania żądania albo wykonywania analizy wystąpi błąd, komunikat zostaje przechwycony i wyświetlony użytkownikowi. Niezależnie od wyniku operacji po zakończeniu żądania stan ładowania jest wyłączany. Dzięki temu interfejs może pokazać użytkownikowi, czy analiza nadal trwa, czy też zakończyła się powodzeniem albo błędem.

W aplikacji klienckiej przygotowano również komponent \texttt{InstructionCard}, którego zadaniem jest przekazanie użytkownikowi podstawowych informacji o sposobie przygotowania nagrania. Wskazówki dotyczą między innymi ustawienia kamery z boku, objęcia w kadrze całej sylwetki i sztangi oraz nagrania pełnej serii ćwiczenia. Ten element nie wpływa bezpośrednio na działanie algorytmu, ale ma znaczenie praktyczne, ponieważ jakość i sposób przygotowania nagrania wpływają na możliwość poprawnego wykrycia punktów kluczowych i późniejszej oceny techniki.

Widok aplikacji klienckiej przed przesłaniem nagrania przedstawiono na rysunku~\ref{fig:widok-poczatkowy-aplikacji}.

\begin{figure}[H]
\centering
\includegraphics[width=0.9\textwidth]{figures/obraz_naglowek.png}
\caption{Widok początkowy aplikacji klienckiej przed przesłaniem nagrania. Źródło: opracowanie własne.}
\label{fig:widok-poczatkowy-aplikacji}
\end{figure}

Na rysunku~\ref{fig:widok-poczatkowy-aplikacji} pokazano formularz przesyłania pliku oraz instrukcje przygotowania nagrania. Elementy te odpowiadają za zebranie danych wejściowych od użytkownika przed rozpoczęciem analizy.

Po wykonaniu analizy otrzymane dane są przekazywane do komponentu \texttt{AnalysisResults}, którego sposób działania opisano w kolejnym podrozdziale.

\subsection{Prezentacja wyników analizy w interfejsie użytkownika}

W momencie zakończenia analizy nagrania przez stronę serwerową aplikacja kliencka otrzymuje odpowiedź i zapisuje ją w stanie komponentu głównego. Następnie otrzymane dane są przekazywane do komponentu \texttt{AnalysisResults}, który odpowiada za zaprezentowanie wyników użytkownikowi. Aplikacja kliencka nie wykonuje dodatkowych obliczeń biomechanicznych, lecz wykorzystuje dane otrzymane z backendu i przekształca je jedynie na potrzeby prezentacji w interfejsie.

Widok wyników został zaprojektowany w taki sposób, aby użytkownik mógł najpierw zobaczyć najważniejsze informacje, a dopiero później przejść do szczegółów. Na początku prezentowana jest nazwa analizowanego pliku, status analizy, liczba wykrytych powtórzeń oraz czas przetwarzania. Pozwala to szybko określić podstawowy wynik działania systemu dla przesłanego nagrania.

Główną częścią widoku jest lista ocen poszczególnych powtórzeń. Każde powtórzenie prezentowane jest osobno za pomocą komponentu \texttt{RepetitionCard}. Taki podział sprawia, że użytkownik nie otrzymuje jednego zbiorczego komunikatu dla całego nagrania, ale może przeanalizować wynik każdego wykrytego powtórzenia oddzielnie. Jest to zgodne z założeniem systemu, ponieważ różnice w wykonaniu mogą występować tylko w części analizowanej serii.

Widok wyników analizy po zakończeniu przetwarzania nagrania przedstawiono na rysunku~\ref{fig:widok-wynikow-analizy}.

\begin{figure}[H]
\centering
\includegraphics[width=0.9\textwidth]{figures/obraz_wyniki.png}
\caption{Widok wyników analizy w aplikacji klienckiej. Źródło: opracowanie własne.}
\label{fig:widok-wynikow-analizy}
\end{figure}

Na rysunku~\ref{fig:widok-wynikow-analizy} widoczny jest główny widok wyników analizy. Wyniki zostały podzielone na osobne karty powtórzeń z odpowiadającymi im statusami.

Każde powtórzenie ma odpowiadającą mu kartę, na której wyświetlany jest jego numer oraz główny status. Do określenia statusu prezentowanego użytkownikowi wykorzystywana jest funkcja pomocnicza \texttt{getRepetitionDisplayStatus}. Dla powtórzeń sklasyfikowanych jako częściowe lub niepoprawne główny status karty odzwierciedla bezpośrednio tę klasyfikację. W pozostałych przypadkach prezentowany jest ogólny status oceny techniki.

Do prezentacji statusów wykorzystano komponent \texttt{StatusBadge}. Jego zadaniem jest przedstawienie statusu danego elementu w powtarzalnej i jednolitej formie wizualnej w różnych częściach interfejsu. Do wyboru odpowiedniej etykiety komponent wykorzystuje funkcję o nazwie \texttt{getStatusLabel}, która zamienia techniczne wartości statusów na bardziej czytelne określenia wyświetlane użytkownikowi, między innymi „Poprawnie”, „Wymaga uwagi”, „Częściowe” i „Niepoprawne”.

Szczegóły oceny pojedynczego powtórzenia nie są widoczne domyślnie. Użytkownik może je rozwinąć za pomocą przycisku dostępnego w karcie powtórzenia. W implementacji odpowiada za to lokalny stan komponentu \texttt{RepetitionCard}, który określa, czy szczegóły są aktualnie pokazane. Takie rozwiązanie ogranicza ilość informacji widocznych na ekranie bezpośrednio po zakończeniu analizy i pozwala użytkownikowi skupić się najpierw na ogólnej ocenie serii.

Po rozwinięciu szczegółów użytkownik może zobaczyć powody wpływające na klasyfikację powtórzenia oraz listę sprawdzeń wykonanych przez system. Powody zapisywane są w polu \texttt{validity\_reasons}, natomiast szczegółowe reguły oceny znajdują się w strukturze \texttt{checks}. Dzięki temu interfejs nie ogranicza się do pokazania samego statusu, ale wskazuje również, które elementy wykonania wpłynęły na wynik.

Za prezentację szczegółowych sprawdzeń odpowiada komponent \texttt{ChecksList}. Dla każdego sprawdzenia wyświetlana jest nazwa ocenianego elementu, odpowiadający mu status oraz komunikat przygotowany przez moduł analizy. Nazwy techniczne reguł są przekształcane na etykiety czytelne dla użytkownika za pomocą funkcji \texttt{formatCheckName}. Dzięki temu techniczne nazwy sprawdzeń są zastępowane bardziej zrozumiałymi określeniami stosowanymi w interfejsie.

Jeżeli system nie wykryje żadnego powtórzenia w przesłanym nagraniu, komponent wyników nie próbuje prezentować pustej listy kart. Zamiast tego wyświetlany jest komunikat informujący, że w filmie nie wykryto powtórzeń. Brak wykrytych powtórzeń jest również możliwym wynikiem działania systemu i powinien zostać jednoznacznie przedstawiony użytkownikowi.

Oprócz oceny pojedynczych powtórzeń w interfejsie prezentowana jest również spójność całej serii. Sekcja ta pojawia się wtedy, gdy wynik analizy zawiera dane zapisane w polu \texttt{consistency\_evaluation}. Wyświetlany jest ogólny status spójności, krótki opis utworzony przez funkcję \texttt{buildConsistencySummary} oraz lista szczegółowych sprawdzeń. Funkcja pomocnicza przekształca informacje otrzymane z modułu analizy w krótszy komunikat przeznaczony do wyświetlenia użytkownikowi.

Przykład rozwiniętych szczegółów oceny pojedynczego powtórzenia oraz spójności serii przedstawiono na rysunku~\ref{fig:szczegoly-powtorzenia}.

\begin{figure}[H]
\centering
\includegraphics[width=0.9\textwidth]{figures/obraz_szczegoly.png}
\caption{Szczegóły oceny powtórzenia oraz spójności serii w interfejsie użytkownika. Źródło: opracowanie własne.}
\label{fig:szczegoly-powtorzenia}
\end{figure}

Na rysunku~\ref{fig:szczegoly-powtorzenia} pokazano sposób prezentacji szczegółowych komunikatów dla pojedynczego powtórzenia. Widok pozwala sprawdzić, które elementy zostały ocenione pozytywnie, a które wymagają dodatkowej uwagi. W dolnej części rysunku przedstawiono sekcję oceny spójności całej serii.

Prezentacja wyników została oparta na podziale informacji na poziom ogólny i szczegółowy. Użytkownik najpierw otrzymuje podstawową ocenę powtórzeń i serii, a następnie może rozwinąć szczegółowe komunikaty dla wybranego powtórzenia.

\subsection{Ograniczenia implementacji przyjętego rozwiązania}

Zaimplementowany system realizuje główne funkcje określone w analizie wymagań. Umożliwia przesłanie nagrania, uruchomienie analizy po stronie serwera, wykrycie powtórzeń, obliczenie wybranych parametrów biomechanicznych oraz prezentację wyników w aplikacji klienckiej. Jednocześnie należy zaznaczyć, że zakres działania systemu został świadomie ograniczony. Wynika to zarówno z charakteru pracy inżynierskiej, jak i z przyjętej metody analizy obrazu wideo.

Pierwsze ograniczenie dotyczy jakości i sposobu przygotowania nagrania. W analizie wymagań założono, że materiał powinien przedstawiać osobę ćwiczącą z boku, obejmować całą sylwetkę oraz pozwalać na wykrycie punktów kluczowych ciała. Implementacja systemu opiera się właśnie na tych założeniach. Jeżeli nagranie jest wykonane z nieodpowiedniej perspektywy, sylwetka jest częściowo zasłonięta albo ćwiczący wychodzi poza kadr, wyniki analizy mogą być niepełne lub mniej wiarygodne. System nie weryfikuje automatycznie wszystkich warunków przygotowania nagrania, dlatego jakość materiału wejściowego ma bezpośredni wpływ na dalszą analizę.

Drugim ograniczeniem jest skupienie się na analizie ruchu w obrazie dwuwymiarowym. W dalszych obliczeniach wykorzystywane są przede wszystkim współrzędne punktów kluczowych w płaszczyźnie obrazu, bez odtwarzania pełnej informacji o głębi sceny. Oznacza to, że wartości kątów i położeń zależą od ustawienia kamery względem osoby ćwiczącej. Wybór strony ciała o wyższej widoczności ogranicza wpływ punktów częściowo zasłoniętych lub słabo widocznych, ale nie usuwa ograniczenia wynikającego z analizy dwuwymiarowej.

Ograniczenie to ma szczególne znaczenie przy ocenie ustawienia pleców i kręgosłupa. W implementacji analizowane są wybrane punkty sylwetki oraz kąty opisujące pozycję tułowia, bioder i kolan. System nie bada jednak rzeczywistej krzywizny kręgosłupa ani ustawienia poszczególnych jego odcinków. Może wskazać sytuacje wymagające uwagi, na przykład niekorzystne ustawienie barku względem biodra, ale nie powinien być traktowany jako narzędzie do diagnozy medycznej. Jest to zgodne z przyjętym założeniem, że system ma wspomagać ocenę techniki, a nie zastępować ocenę specjalisty.

Kolejne ograniczenie wynika z użycia gotowego modelu estymacji pozy człowieka. Zwracane przez model punkty kluczowe są przewidywanymi położeniami, a nie dokładnymi pomiarami wykonanymi bezpośrednio na ciele osoby. Jakość detekcji może zależeć między innymi od oświetlenia, ubioru, rozdzielczości obrazu oraz obecności innych obiektów w kadrze. Błędy detekcji mogą wpływać na wartości obliczanych parametrów oraz końcową ocenę techniki.

Z tym ograniczeniem wiąże się również sposób wyznaczania położenia gryfu sztangi. W obecnej implementacji położenie sztangi jest przybliżane jako środek pomiędzy lewym i prawym nadgarstkiem. System nie wykorzystuje osobnego modelu detekcji gryfu ani talerzy obciążenia. Rozwiązanie pozwala oszacować pionowy przebieg ruchu sztangi, jednak jego dokładność zależy od poprawności detekcji nadgarstków. Problem może być szczególnie widoczny w sytuacjach, w których dłonie są częściowo zasłonięte albo model estymacji pozy nie wykrywa ich stabilnie.

Ocena techniki została zaimplementowana jako zestaw reguł decyzyjnych, a nie jako osobny model uczenia maszynowego nauczony oceny techniki na podstawie zbioru danych. Ocena ogranicza się więc do parametrów i reguł jawnie zaimplementowanych w systemie. Nie uwzględnia wszystkich możliwych różnic anatomicznych, indywidualnej mobilności ani innych cech wykonania, które nie zostały opisane przez przyjęte metryki.

System nie dokonuje również analizy sił działających na ciało osoby ćwiczącej ani momentów działających w stawach. Ocena opiera się na geometrii ruchu wyznaczonej na podstawie obrazu, a nie na bezpośrednich pomiarach wielkości dynamicznych, takich jak siły reakcji podłoża czy obciążenia działające na poszczególne stawy.

Ograniczenia występują również po stronie aplikacyjnej. W obecnej wersji system analizuje pojedyncze nagranie przesłane przez użytkownika. Nie zaimplementowano kont użytkowników, historii treningów, porównywania wyników między sesjami ani generowania planów treningowych. Zakres ten jest zgodny z wymaganiami funkcjonalnymi, w których skupiono się na analizie pojedynczego materiału wideo oraz prezentacji wyniku.

Analiza wykonywana jest synchronicznie w ramach obsługi żądania HTTP, dlatego czas obsługi żądania zależy od czasu potrzebnego na przetworzenie nagrania. Rozwiązanie to może stanowić ograniczenie w przypadku dłuższych materiałów lub równoczesnej obsługi większej liczby analiz. W dalszym rozwoju projektu możliwe byłoby przeniesienie analizy do kolejki zadań lub osobnego procesu roboczego.

Ostatnim ograniczeniem jest zakres weryfikacji skuteczności reguł oceny. System został przygotowany tak, aby umożliwić analizę nagrania i przedstawić wynik w uporządkowanej formie, jednak szersza walidacja jakości ocen wymagałaby wykorzystania większego zbioru nagrań oraz porównania wyników systemu z oceną ekspertów.

Opisane ograniczenia wyznaczają zakres, w jakim należy interpretować wyniki systemu, oraz wskazują obszary możliwe do rozwinięcia w kolejnych wersjach projektu.
